\section{Đánh giá}
\label{sec:evaluation}

\subsection{Phương pháp kiểm tra}

Đánh giá được thực hiện trên ba lớp: kiểm thử code local, replay dữ liệu công khai trên một ngày Sandbox, và lặp lại replay trong Docker/WSL. Kiểm thử dùng các đầu ra cụ thể thay vì chỉ kiểm tra service có khởi động được hay không.

\begin{xltabular}{\textwidth}{>{\raggedright\arraybackslash}p{4.4cm} >{\raggedright\arraybackslash}X}
\caption{Kết quả kiểm chứng của implementation}
\label{tab:evaluation-results} \\
\hline
\textbf{Kiểm chứng} & \textbf{Kết quả quan sát được} \\
\hline
\endfirsthead
\hline
\textbf{Kiểm chứng} & \textbf{Kết quả quan sát được} \\
\hline
\endhead
\hline
\endfoot
\hline
\endlastfoot

Unit test và simulation
& 153 test Pytest pass. Fixture 9 dòng được định tuyến thành 4 valid, 2 dead-letter và 4 quality rows; cũng xuất được CSV quality log. \\

Replay dữ liệu công khai
& 100/100 event được publish từ ngày 01/11/2020 qua emulator và DirectRunner, ghi thành 100 dòng JSONL hợp lệ. \\

Kho trên Sandbox
& \texttt{events\_20201101} chứa 100 event. dbt tạo 19/19 model và pass 5/5 test. \texttt{fct\_events}=100, \texttt{fct\_sessions}=17; ba KPI view có thể query và mỗi view trả 11 dòng. \\

Dashboard thực tế
& Streamlit đọc được phiên/SKU từ BigQuery thật và render 0 exception. Health endpoint và trang chính trả HTTP 200. \\

Docker offline
& 4 valid event, 2 dead-letter, 4 quality rows từ fixture. Emulator và pipeline được dừng, container được xóa. \\

Docker Sandbox
& 100 event được load vào dataset riêng \texttt{marketing\_dss\_docker\_smoke}. Kết quả được 19/19 model, 5/5 dbt test và ba KPI view, với \texttt{fct\_events}=100 và \texttt{fct\_sessions}=17. \\

Benchmark giới hạn
& Query raw đọc 7.68 MiB trong 1.636 giây; bốn query KPI lần lượt tốn 0.787, 0.787, 0.827 và 0.857 giây. Trình chạy tắt query cache ở query thật để đo latency. \\
\end{xltabular}

\subsection{Quota và tài nguyên}

Quá trình phát lại (replay) 100 sự kiện đầy đủ từ nguồn dữ liệu trong ngày kiểm thử tiêu tốn 22.47 MiB trên BigQuery. Con số này lớn hơn kích thước thực tế của các thông điệp do truy vấn cần quét toàn bộ các trường hợp lệ để dựng bản ghi JSON. Biến cấu hình \texttt{BQ\_MAX\_QUERY\_BYTES=134217728} giới hạn dung lượng quét của mỗi truy vấn ở mức tối đa 128 MiB: cả publisher và dashboard đều chạy kiểm tra trước (dry-run) lẫn áp trần khi thực thi, trong khi dbt cũng được cấu hình ngưỡng giới hạn tương tự trong profile. Ngoài ra, số liệu sử dụng hằng tháng được trích xuất trực tiếp từ \texttt{INFORMATION\_SCHEMA.JOBS\_BY\_PROJECT} thay vì chỉ ước tính qua tệp nhật ký nội bộ. \medskip

Tính đến thời điểm kết thúc đợt kiểm thử, toàn bộ dự án đã xử lý tổng cộng \textbf{103,485,196 bytes} (khoảng 103.49 MB), tương đương 0.19\% hạn mức ngân sách 50 GiB tự đặt. Lưu ý rằng nhật ký lịch sử job của hệ thống có thể cập nhật trễ và số liệu này chỉ mang tính theo dõi chứ không phải hóa đơn thanh toán thực tế. Việc thiết lập trần ngân sách đóng vai trò như một chốt chặn để rà soát trước khi mở rộng quy mô dữ liệu, chứ không bảo đảm mọi truy vấn trong tương lai đều tiêu tốn ít tài nguyên. \medskip

Khi chạy kiểm thử replay trên môi trường Docker, mức tiêu thụ RAM ghi nhận được lần lượt là khoảng 211 MiB đối với emulator, 212 MiB cho pipeline và 86.49 MiB cho dashboard. Tiến trình watchdog giám sát liên tục mức sử dụng RAM, dung lượng ổ đĩa và các PID; khi kết thúc, các tiến trình con trên Windows được thu hồi theo cây tiến trình, còn trên Linux sử dụng cơ chế process group cho emulator. Sau khi hoàn tất phiên kiểm thử, hệ thống đã dọn dẹp sạch tài nguyên, không còn tiến trình nào chiếm dụng cổng 8085 hay 8501 cũng như không còn container Compose tồn đọng.

\subsection{Những giới hạn của đánh giá}

Thứ nhất, dù tập dữ liệu nguồn GA4 kéo dài 92 ngày, kịch bản kiểm thử replay mới chỉ chạy trên 100 sự kiện đầu tiên của một ngày. Kết quả này chưa đủ cơ sở để đánh giá hiệu năng hay mức tiêu thụ quota khi xử lý toàn bộ 92 ngày dữ liệu. Bên cạnh đó, các thông số độ trễ nêu trên được đo tuần tự trên tập mẫu nhỏ nên không thể xem là cam kết hiệu năng cho các tập dữ liệu quy mô lớn. \medskip

Thứ hai, lô dữ liệu thử nghiệm không phát sinh sự kiện mua hàng hay dữ liệu mặt hàng. Vì vậy, đợt kiểm chứng này chủ yếu xác nhận tính đúng đắn của luồng thu thập, chuẩn hóa lược đồ và gom nhóm phiên. Dù các quy tắc xử lý giao dịch và sản phẩm vẫn được kiểm tra qua bộ dữ liệu giả lập (fixture) cục bộ, hệ thống chưa ghi nhận được giao dịch có doanh thu thực tế có giá trị dương trên môi trường Sandbox. Để các chỉ số KPI doanh thu mang lại giá trị phân tích thực tiễn, cần mở rộng quy mô lô dữ liệu thử nghiệm đi kèm bước ước tính chi phí quét (dry-run) và cấp phát quota riêng. \medskip

Cuối cùng, DirectRunner và Pub/Sub Emulator vốn chỉ phục vụ mục đích kiểm thử trên máy cục bộ, nên không phản ánh được khả năng lưu trữ bền vững, cơ chế phân quyền IAM hay hiệu năng thực tế của các dịch vụ phân tán trên đám mây. Tương tự, tệp xuất CSV cho Looker Studio chỉ là bản chụp dữ liệu tĩnh (snapshot) phục vụ minh họa, chưa phải kết nối trực tiếp (live) tới BigQuery. Các phương án nâng cấp hạ tầng và kết nối dữ liệu tự động này được tách riêng sang phần hướng mở rộng thay vì tính vào phạm vi đã hoàn thiện. \medskip
